\documentclass[11pt]{article}
\usepackage[a4paper,margin=27mm]{geometry}
\usepackage[T1]{fontenc}
\usepackage{lmodern,microtype,amsmath,amssymb,amsthm,mathtools,booktabs}
\usepackage[colorlinks=true,linkcolor=blue,urlcolor=blue,citecolor=blue]{hyperref}
\usepackage{xurl}
\newtheorem{theorem}{Theorem}[section]
\newtheorem{lemma}[theorem]{Lemma}
\newcommand{\ii}{\mathrm{i}}
\newcommand{\Li}{\operatorname{Li}}
\newcommand{\Log}{\operatorname{Log}}
\newcommand{\dd}{\,\mathrm{d}}
\newcommand{\E}{\mathcal{G}}
\title{Asymptotics and inverses for a stable Hilbert series\\\large The coefficients of OEIS A126348}
\author{}
\date{1 October 2026}
\begin{document}
\maketitle
\begin{abstract}
We derive a relative all-orders expansion for the coefficients of
$\prod_{k\ge1}(1+q^k/(1-q))$, a stable Hilbert series associated with top-degree
Grothendieck and Lascoux polynomials. A quadratic logarithmic saddle provides
a resummed leading term and rational correction functions, with a finite rule
for every order and a uniform remainder. An exact Jacobi and eta factorization
separately identifies convergent oscillatory sectors of the generating function.
The coefficient theorem yields an explicit inverse with $o(1)$ error in the
index on the sequence range, and qualified integer threshold rules.
\end{abstract}

\section{The coefficient theorem}
Define
\begin{equation}\label{eq:F}
 F(q)=\prod_{k\ge1}\left(1+\frac{q^k}{1-q}\right)
     =\sum_{n\ge0}a_nq^n,\qquad |q|<1.
\end{equation}
The sequence begins $1,1,2,4,7,12,20,33,53,84,131$ and is
A126348~\cite{oeis}. It counts colored partitions with distinct
colors~\cite{NR}. Pan and Yu identify \eqref{eq:F} as the Hilbert series
of the stable span of top-degree Grothendieck and Lascoux
polynomials~\cite[Proposition 6.10]{PY}.

Put $c=\pi^2/6$ and
\[
 A(L)=\tfrac12L^2+L+c,\quad B(L)=L^2+L+2c,\quad
 V(L)=L^2+3L+2c+1.
\]
For sufficiently large $n$, let $L=L_n>0$ be the unique solution of
\begin{equation}\label{eq:saddle}
 n=e^{2L}A(L),\qquad t=e^{-L}.
\end{equation}
All expansion orders are fixed while $n\to\infty$.

\begin{theorem}\label{thm:coeff}
For each integer $J\ge0$, there are rational functions
$C_j\in\mathbb{Q}(c,L)$, with $C_0=1$, such that
\begin{equation}\label{eq:coeff}
 a_n=\frac{\exp\{e^LB(L)-1\}}{\sqrt{2\pi e^{3L}V(L)}}
 \left(\sum_{j=0}^{J}e^{-jL}C_j(L)
 +O_J\bigl(e^{-(J+1)L}(1+L)^{J+1}\bigr)\right).
\end{equation}
Moreover, $C_j(L)=O_j((1+L)^j)$. A finite exact rule for all $C_j$ is
provided in Section~\ref{sec:rule}.
\end{theorem}
This is a coefficient equivalent, not only a logarithmic equivalent.
Keeping \eqref{eq:saddle} exact retains infinitely many corrections in
inverse powers of $\log n$.

For explicit formulas define
\begin{equation}\label{eq:Vrec}
 V_0=\tfrac12L^2+c,\qquad V_j=(j+\partial_L)V_{j-1}.
\end{equation}
Thus $V_1=A$, $V_2=V$, and
\begin{align*}
 V_3&=3L^2+11L+6c+6,\\
 V_4&=12L^2+50L+24c+35,\\
 V_5&=60L^2+274L+120c+225,\\
 V_6&=360L^2+1764L+720c+1624.
\end{align*}
Let $P_1=(5-L)/24$, $P_2=-1/9$, and
\begin{align}
 E_1&=\frac{V_4}{8V^2}-\frac{5V_3^2}{24V^3},\label{eq:E1}\\
 E_2&=-\frac{V_6}{48V^3}+\frac{7V_3V_5}{48V^4}
 +\frac{35V_4^2}{384V^4}-\frac{35V_3^2V_4}{64V^5}
 +\frac{385V_3^4}{1152V^6}.\label{eq:E2}
\end{align}
The first two corrections are
\begin{align}
 C_1&=P_1+E_1,\label{eq:C1}\\
 C_2&=P_2+\tfrac12P_1^2+P_1E_1+E_2
       -\frac1{48V}-\frac{(6-L)V_3}{48V^2}.\label{eq:C2}
\end{align}
The last two terms of $C_2$ arise from the displaced amplitude.

\section{An exact modular factorization}\label{sec:modular}
For real $t>0$ set
\[
 q=e^{-t},\quad \delta=1-q,\quad \lambda=-\log\delta,\quad
 u=\lambda/t-1/2,\quad Q=e^{-4\pi^2/t}.
\]
Use $(z;q)_\infty=\prod_{k\ge0}(1-zq^k)$.
\begin{theorem}\label{thm:modular}
The following identity is exact:
\begin{align}
 F(e^{-t})&=e^{K(t)}\frac{\Theta(u,t)}{(Q;Q)_\infty},\label{eq:modular}\\
 K(t)&=\frac{\lambda^2/2+c}{t}-\frac\lambda2+\frac t{12}-G(t),\label{eq:K}\\
 G(t)&=\log(-\delta;e^{-t})_\infty
 =\sum_{j\ge1}\frac{(-1)^{j+1}\delta^j}{j(1-e^{-jt})},\label{eq:G}\\
 \Theta(u,t)&=\sum_{m\in\mathbb Z}e^{-2\pi^2m^2/t}e^{2\pi\ii m u}.
\end{align}
Writing $p(k)$ for the ordinary partition function, with $p(0)=1$, the
normalized quotient has the exact convergent expansion
\begin{align}
 \frac{\Theta(u,t)}{(Q;Q)_\infty}
 &=\sum_{r\ge0}D_r(u)e^{-2\pi^2r/t},\label{eq:sectors}\\
 D_r(u)&=\sum_{\substack{m\in\mathbb Z\,,\ m^2\le r\\r-m^2\ {\rm even}}}
 p\left(\frac{r-m^2}{2}\right)e^{2\pi\ii m u}.\label{eq:Dr}
\end{align}
For fixed $R$, the remainder after $r=R$ is
$O_R(e^{-2\pi^2(R+1)/t})$ as $t\downarrow0$, uniformly in real $u$.
\end{theorem}
\begin{proof}
Jacobi's triple product~\cite[17.8.1]{DLMF} gives
\[
 F(q)(-\delta;q)_\infty(q;q)_\infty
 =\sum_{k\in\mathbb Z}e^{-tk^2/2+(\lambda-t/2)k}.
\]
Completing the square and applying Poisson summation turns the right side into
\[
 \sqrt{2\pi/t}\,e^{(\lambda-t/2)^2/(2t)}\Theta(u,t).
\]
The eta transformation~\cite[17.2.6\_1]{DLMF} is
\[
 (q;q)_\infty=\sqrt{2\pi/t}\,e^{-\pi^2/(6t)+t/24}(Q;Q)_\infty.
\]
The square-root factors cancel and $t/8-t/24=t/12$.
Expanding the logarithms in the partner product gives \eqref{eq:G},
absolutely convergently since $0<\delta<1$. Multiply the theta series by
$(Q;Q)_\infty^{-1}=\sum_{k\ge0}p(k)Q^k$ to obtain \eqref{eq:Dr}.
Its convergence and tail follow from $p(k)=e^{O(\sqrt k)}$.
\end{proof}
The first amplitudes are
\[
 D_0=1,\quad D_1=2\cos(2\pi u),\quad D_2=1,\quad
 D_3=2\cos(2\pi u),\quad D_4=2+2\cos(4\pi u).
\]
These are exactly normalized oscillatory sectors of the generating
function relative to the exact core $e^{K(t)}$. They do not by themselves
identify exponentially small coefficient sectors.

\section{The uniform logarithmic expansion}\label{sec:log}
Let $L=\log(1/t)$. Expansion of the exact core gives
\begin{equation}\label{eq:logF}
 \log F(e^{-t})=\frac{L^2/2+c}{t}-1+
 \sum_{j=1}^{J}t^jP_j(L)+O_J(t^{J+1}(1+L)),
\end{equation}
where every $P_j$ is affine in $L$. In addition to $P_1,P_2$ above,
\[
 P_3=\frac{2L+245}{5760},\quad P_4=-\frac{11}{900},\quad
 P_5=-\frac{8L+2037}{1451520},\quad P_6=\frac{341}{52920}.
\]
For a finite coefficient rule, rewrite the $j$th term of $G$ as
\[
 \frac{(-1)^{j+1}(1-e^{-t})^{j-1}}
 {j(1+e^{-t}+\cdots+e^{-(j-1)t})}.
\]
Only $j\le J+1$ contributes through degree $J$. Also
$\lambda=L+t/2-t^2/24+t^4/2880-\cdots$.
In \eqref{eq:K}, the order-zero terms $L/2$ and $-L/2$ cancel, and
the first term of $G$ is exactly one.

For an independent symbolic route, Euler--Maclaurin gives
\[
 \log F(e^{-t})\sim-\frac{\Li_2(-e^\lambda)}{t}
 -\frac12\log(1+e^\lambda)
 -\sum_{m\ge1}\frac{B_{2m}t^{2m-1}}{(2m)!}\Li_{2-2m}(-e^\lambda).
\]
The reflection identity
$-\Li_2(-e^\lambda)=\lambda^2/2+c+\Li_2(-e^{-\lambda})$
recovers the same coefficients. We use the exact modular identity to
justify the complex remainder.

\begin{lemma}\label{lem:complex}
For a fixed sufficiently large $K_0$, uniformly on
$|z-t|\le t/(K_0L^2)$, formula \eqref{eq:logF} holds with $t$ replaced by
$z$ and its displayed $L$ replaced by $\Log(1/z)$. The remainder is still
$O_J(t^{J+1}(1+L))$.
\end{lemma}
\begin{proof}
For small $t$, $|\delta(z)|\le2t<1$ and $\Re z\ge t/2$, so
\[
 |1-e^{-jz}|\ge1-e^{-j\Re z}\ge c_0\min(jt,1).
\]
The tail of \eqref{eq:G} beyond $j=J+1$ is bounded by
$C\sum_{j\ge J+2}(2t)^j/[j\min(jt,1)]=O_J(t^{J+1})$.
Each earlier term has a removable singularity at zero. The function
$\lambda(z)-\Log(1/z)$ is analytic there as well; the explicit logarithm
accounts for the factor $1+L$.

Furthermore $\Re(1/z)\sim1/t$ and
$|\Im(\lambda(z)/z)|=O(1/(tL))$. The nonzero theta terms are bounded by
\[
 \exp\{-c_1m^2/t+C_1|m|/(tL)\},
\]
whose sum is $O(e^{-c_2/t})$. The dual Euler product obeys a similar
bound. Analytic continuation of \eqref{eq:modular} is valid on this
neighborhood: $\delta(z)$ has positive real part, the small-fugacity
partner has no zero, and the theta factor is close to one. These
observations also specify the logarithm branches.
\end{proof}
The convergent exact $j$-series is not a convergence assertion for its
power-log expansion at zero. Its individual terms have boundary
root-of-unity singularities accumulating there. Our power-log statements
are finite-truncation asymptotic statements.

\section{Coefficient extraction with the precise remainder}\label{sec:proof}
Write $r=e^{-t}$, $q=re^{\ii\theta}$, $D=|1-q|$, and
$\alpha=(1-r)/D$. The triangle inequality gives
$|F(q)|\le\prod_{k\ge1}(1+r^k/D)$.
For $k\le\lfloor\lambda/t\rfloor$, $w_k=r^k/(1-r)\ge1$, and
\[
 \log(1+w_k)-\log(1+\alpha w_k)\ge(1-\alpha)/2.
\]
Since $D^2=(1-r)^2+2r(1-\cos\theta)$, uniformly for $|\theta|\le\pi$,
\begin{equation}\label{eq:minor}
 \frac{|F(re^{\ii\theta})|}{F(r)}
 \le\exp\left\{-c_3\frac Lt
       \min\left(\frac{\theta^2}{t^2},1\right)\right\}.
\end{equation}
At $|\theta|\ge t/(K_0L^2)$ this is
$\exp\{-\Omega(1/(tL^3))\}$, negligible to every power of $t$.
All other roots of unity are suppressed simultaneously.

Put $S(t)=(L^2/2+c)/t$. Then $-S'(t)=n$ and
$(-1)^kS^{(k)}(t)=t^{-k-1}V_k(L)$.
Cauchy's formula reads
\[
 a_n=\frac{e^{nt}}{2\pi}\int_{-\pi}^{\pi}
 F(e^{-t+\ii\theta})e^{-\ii n\theta}\dd\theta.
\]
Set $\epsilon=\sqrt t$ and $X=\theta\sqrt V/t^{3/2}$, with central
cutoff $|X|\le\epsilon^{-1/4}$. On the intermediate arc up to
$|\theta|=t/(K_0L^2)$, Taylor's theorem gives
\[
 \Re\{S(t-\ii\theta)-S(t)-\ii n\theta\}\le-c_4X^2.
\]
The cubic-to-quadratic ratio is $O(|\theta|/t)=O(L^{-2})$.
The remaining logarithmic amplitude changes by $O(t(1+L))$, which
cannot affect this bound beyond the central cutoff.

On the central arc, Lemma~\ref{lem:complex} replaces the amplitude by its
finite sum through degree $J$, with multiplicative error
$O_J(t^{J+1}(1+L))$. The normalized integrand is
$e^{-X^2/2}\exp(Q+R_J)$ with $Q,R_J$ as in Section~\ref{sec:rule}.
For a precise remainder, treat $\epsilon$ as complex while $L$ and real
$X$ are fixed. On the auxiliary circle
\[
 |\epsilon|=\rho=\frac{c_J}{(1+L)^{1/2}(1+|X|)^3},
\]
$Q$ and $R_J$ are uniformly bounded. Indeed,
$|w|=|\epsilon X|/\sqrt V\ll1$,
$Q=O(|\epsilon|\,|X|^3/(1+L))$, and
$R_J=O_J(|\epsilon|^2(1+L))$.
The exact phase has a removable singularity at $\epsilon=0$ after its
linear and quadratic terms have been separated. Cauchy's Taylor
remainder through degree $2J+1$ is therefore bounded by
\[
 C_J|\epsilon|^{2J+2}(1+L)^{J+1}(1+|X|)^{6J+6}.
\]
For actual $\epsilon=\sqrt t$, $|\epsilon|\le\rho/2$ on the central
cutoff for large $L$, because $\epsilon^{1/4}(1+L)^{1/2}\to0$.
Gaussian integration yields exactly
$O_J(t^{J+1}(1+L)^{J+1})$.
Odd Taylor terms vanish on the symmetric interval. Extending the
finitely many even moments to the real line costs an exponentially
small error. Expansion through the odd degree before cancellation
avoids a half-power loss.

The Jacobian is $t^{3/2}/\sqrt V$ and $S(t)+nt=t^{-1}B(L)$.
The constant amplitude is $e^{-1}$. These facts prove
Theorem~\ref{thm:coeff}, including its coefficient growth bound.

\section{A finite coefficient rule}\label{sec:rule}
Let $\E$ be the Gaussian moment functional
$\E(X^{2m})=(2m-1)!!$, $\E(X^{2m+1})=0$.
For $w=-\ii\epsilon X/\sqrt V$, set
\begin{align*}
 Q&=\sum_{k\ge3}\frac{\ii^k\epsilon^{k-2}V_kX^k}{k!V^{k/2}},\\
 R_J&=\sum_{r=1}^{J}\epsilon^{2r}(1+w)^rP_r(L-\log(1+w)).
\end{align*}
Then for $0\le j\le J$,
\begin{equation}\label{eq:rule}
 C_j(L)=\E\left([\epsilon^{2j}]\exp(Q+R_J)\right).
\end{equation}
Only $k\le2j+2$ and $r\le j$ enter. Gaussian parity removes half-powers
of $t$ and square roots of $V$, giving the asserted rational functions.
For implementation, if $H(\epsilon)=\sum_{k\ge1}h_k\epsilon^k$ and
$e^H=\sum b_k\epsilon^k$, use
\[
 b_0=1,\qquad b_n=\frac1n\sum_{k=1}^{n}k h_kb_{n-k}.
\]
The included independent symbolic reconstruction verifies
\eqref{eq:C1}--\eqref{eq:C2}, odd cancellation, and computes $C_3$.

\section{Growth inversion}\label{sec:inverse}
For a large target $Y$, put $y=\log Y$ and choose a new $L$ from
\begin{equation}\label{eq:carrier}
 y=H(L):=e^LB(L).
\end{equation}
This inverse-carrier coordinate differs from the forward saddle at a
preassigned $n$. Set
\[
 N_0=e^{2L}A(L),\quad T=1+\tfrac32L+\tfrac12\log(2\pi V),\quad
 T'=\tfrac32+V'/(2V).
\]
\begin{theorem}\label{thm:inverse}
The explicit function
\begin{equation}\label{eq:Nhat}
 \widehat N(Y)=N_0+e^LT+\frac{T^2+2TT'}{2V}-C_1(L)
\end{equation}
satisfies
$\widehat N(a_n)=n+O(e^{-L}(1+L)^3)=n+o(1)$.
Nearest-integer rounding therefore recovers the index on the sequence
range for every sufficiently large $n$.
\end{theorem}
\begin{proof}
Write $N(L)=e^{2L}A(L)$. Exactly,
\begin{equation}\label{eq:derivs}
 H'=e^LV,\qquad N'=e^{2L}V,\qquad \frac{\dd N}{\dd H}=e^L.
\end{equation}
Define a positive smooth first-correction model $M_1(x)$ by the leading
factor in \eqref{eq:coeff} times $1+tC_1$, where $x=N(L)$.
It is eventually increasing since $(\log M_1)'(x)=t(1+o(1))$.
The forward logarithmic error $O(t^2L^2)$ becomes $O(tL^2)$ under its
inverse.

For a direct residual verification, let $H(L)=y$ and set
\begin{align*}
 \Delta_1&=tT/V,\\
 \Delta_2&=t^2\left(\frac{TT'}{V^2}
 -\frac{(V+V')T^2}{2V^3}-\frac{C_1}{V}\right).
\end{align*}
Substituting $\Delta_1+\Delta_2$ into
\[
 H(L+\Delta)-T(L+\Delta)
 +\log\{1+e^{-(L+\Delta)}C_1(L+\Delta)\}=y
\]
leaves residual $O(t^2(1+L)^2)$. Its derivative is asymptotic to
$t^{-1}V$, so the root error is $O(t^3)$. Expanding $N(L+\Delta)$ through
degree two gives exactly \eqref{eq:Nhat}, with index error
$O(t(1+L)^2)$, proving the stated weaker bound.
\end{proof}
The term $-C_1(L)$ is unbounded of order $L$ and cannot be discarded in
an inverse intended to recover integer indices.

For all orders let
\[
 h(L)=-T+\log\left(1+\sum_{j\ge1}e^{-jL}C_j(L)\right),\qquad
 D_y=e^{-L}V^{-1}\partial_L.
\]
The formal reversion rule is
\begin{equation}\label{eq:allinverse}
 N(L)+\sum_{k\ge1}\frac{(-1)^k}{k!}
 D_y^{k-1}\{e^Lh(L)^k\},\qquad H(L)=y.
\end{equation}
At each exponential order it involves finitely many terms. A finite
smooth forward model and the same Taylor-residual proof turn each
sufficiently deep truncation into a controlled range inverse.
For $D_2=C_2-C_1^2/2$, the next correction after \eqref{eq:Nhat} is
$e^{-L}B_1$, where
\[
 B_1=-D_2-\frac{(TC_1)'}{V}
 +\frac1{6V}\left(\frac{T^3+3T^2T'}{V}\right)'.
\]
A Lambert starting value for \eqref{eq:carrier} is
$L_0=2W(\sqrt y/2)$, solving $y=e^{L_0}L_0^2$; inverse-$L_0$
corrections account for the rest of $B$. Keeping the carrier exact is
usually simpler. In particular,
\[
 \widehat N(Y)\sim\frac{(\log Y)^2}{2(\log\log Y)^2}.
\]

\section{Thresholds and scope}
The $k=1$ factor of $F$ is $(1-q)^{-1}$. The remaining product has
nonnegative coefficients, and its $k=2$ factor has positive coefficients
in every degree at least two. Thus $a_n>a_{n-1}$ for all $n\ge2$.
The threshold inverse $N(Y)=\min\{n\ge2:a_n\ge Y\}$ is well defined.
For sufficiently large $Y$, it lies between the ceilings of
\[
 \widehat N(Y)\pm Ce^{-L}(1+L)^3
\]
for an appropriate constant $C$. This follows from monotonicity of the
smooth model and the uniform forward error at neighboring integers.
Eventually there are at most two candidate ceilings. A single ceiling
is justified only when both bracket endpoints have the same ceiling.
At $Y=a_n$, use nearest-integer rounding instead.

Formal inversion, inversion on the sequence range, a threshold staircase,
and inversion of an exact continuous interpolation are different objects.
This paper treats the first three. It specifies no exact sequence
interpolation, effective onset, coefficient-sector amplitudes, or Stokes
classification. The normalized modular sectors belong to the generating
function itself.

\section{Checks and literature scope}
The accompanying scripts reconstruct $P_1$ through $P_6$ independently
from Euler--Maclaurin and the exact small-fugacity series. Gaussian
extraction independently reproduces $C_1,C_2$ and odd cancellation.
Exact integer coefficient generation and high-precision evaluation give
these selected inverse errors:
\begin{center}
\begin{tabular}{rrr}\toprule
$n$ & $N_0(a_n)-n$ & $\widehat N(a_n)-n$\\\midrule
100 & $-24.644779$ & $-0.050465$\\
1000 & $-81.976080$ & $-0.015455$\\
10000 & $-261.371189$ & $-0.005225$\\\bottomrule
\end{tabular}
\end{center}
The modular identity was checked to at least ninety decimal places at
five positive arguments. These are diagnostics, not an effective-onset
certificate; the theorems rest on the analytic proof.

The live OEIS entry inspected on 1 October 2026 contains no asymptotic
formula. The cited combinatorial papers establish interpretations and
generating functions. Classical q-product expansions~\cite{McI} and a
recent uniform q-Pochhammer expansion~\cite{AR} are relevant background.
The latter's stated uniform corollary imposes a fixed lower bound on the
real part of its logarithmic argument, which our moving fugacity leaves.
The exact pairing here avoids that issue. A bounded literature search
did not locate a published coefficient or inverse-growth theorem for
A126348; this does not assert exhaustive absence.

\begin{thebibliography}{9}
\bibitem{oeis} OEIS Foundation, A126348. \url{https://oeis.org/A126348}.
\bibitem{PY} J. Pan and T. Yu, Top-degree components of Grothendieck and
Lascoux polynomials, \emph{Algebraic Combinatorics} \textbf{7}(1)
(2024), 109--135. \url{https://doi.org/10.5802/alco.326}.
\bibitem{NR} B. J. Naranjo and J. L. Ram\'irez, On colored partitions and
Euler-type identities, \emph{Integers} \textbf{26} (2026), A20.
\url{https://math.colgate.edu/~integers/aa20/aa20.pdf}.
\bibitem{DLMF} NIST Digital Library of Mathematical Functions, Sections
17.2 and 17.8. \url{https://dlmf.nist.gov/17.2}; \url{https://dlmf.nist.gov/17.8}.
\bibitem{McI} R. J. McIntosh, Some asymptotic formulae for q-shifted
factorials, \emph{Ramanujan Journal} \textbf{3} (1999), 205--214.
\url{https://doi.org/10.1023/A:1006949508631}.
\bibitem{AR} A. Arabi Ardehali and H. Rosengren, A new product formula for
$(z;q)_\infty$, with applications to asymptotics,
\emph{Constructive Approximation} (2026), published 19 September.
\url{https://doi.org/10.1007/s00365-026-09783-2}.
\end{thebibliography}
\end{document}
